\section*{Results}

\subsection*{The model on a random split, and what the random split measures}

On the random split the full stack reached a mean average precision of 0.850 over the seven drugs
(Table~\ref{tab:main}). The average hides the interesting part. Ciprofloxacin reached 0.975 and
ampicillin 0.943, while gentamicin reached 0.709 and amoxicillin--clavulanic acid 0.754: a spread
of 0.27 between the best and worst drug for one architecture trained on one set of genomes. The
accuracy column moves the other way. Gentamicin, the weakest drug by average precision, has an
accuracy of 0.936, because only 12\% of its strains are resistant and a model can be accurate while
missing most of them. That is the trap the metric choice is meant to avoid, and it recurs in
clinical form below.

\begin{table}[t]\centering\small
\caption{\textbf{The full stack on the random split}, per drug, at a fixed threshold of 0.5.
AP, average precision; MCC, Matthews correlation coefficient; CA, categorical agreement, which
equals accuracy here because intermediate results were collapsed into resistant.}
\label{tab:main}
\begin{tabular}{lrrrrrrr}
\toprule
Drug & Prevalence & AP & 95\% CI & MCC & $F_1$ & $F_2$ & CA \\
\midrule
CIP & 0.228 & 0.975 & [0.957, 0.989] & 0.919 & 0.936 & 0.932 & 0.973 \\
AMP & 0.501 & 0.943 & [0.919, 0.963] & 0.795 & 0.897 & 0.895 & 0.897 \\
CTX & 0.152 & 0.904 & [0.853, 0.945] & 0.790 & 0.817 & 0.813 & 0.953 \\
SXT & 0.310 & 0.842 & [0.779, 0.899] & 0.702 & 0.794 & 0.787 & 0.873 \\
CAZ & 0.131 & 0.823 & [0.754, 0.877] & 0.689 & 0.726 & 0.716 & 0.934 \\
AMC & 0.312 & 0.754 & [0.687, 0.813] & 0.528 & 0.658 & 0.626 & 0.805 \\
GEN & 0.118 & 0.709 & [0.619, 0.790] & 0.626 & 0.655 & 0.611 & 0.936 \\
\midrule
Mean & --- & 0.850 & --- & 0.721 & 0.783 & 0.768 & 0.910 \\
\bottomrule
\end{tabular}
\end{table}

Three checks bound what that 0.850 means. Near-duplicate strains are not driving it: removing the 47
test strains within 0.5\% SNP divergence of a training strain changed the mean by $-0.005$, with the
largest single drop 0.023 (CAZ) and one drug improving (AMC, $+0.008$). The intervals in
Table~\ref{tab:main} are wide enough that no comparison between closely related variants of the
model clears them; gentamicin alone spans 0.17. And every model here was trained once under a fixed
seed, so these intervals cover which strains fall in the test partition, not training variability.

The third check is the one that reframes the paper. A model that never read a base of sequence,
using only the strain's country, year, submitting study, sequencing platform and five assembly
statistics, reached a mean average precision of 0.705 on the same split --- 83\% of the full
genomic model's score (Figure~\ref{fig:floor}). Assembly statistics alone reached 0.601 and the
epidemiological fields alone 0.516. None of this is resistance biology, and none of it survives a
split that moves populations: across the two country splits and the temporal split within Norway,
the epidemiological fields scored an \nap{} between 0.000 and 0.056, and the submitting study alone
scored exactly 0.000 on three of those four splits. Assembly statistics held an \nap{} of
0.097--0.232 there, and on the lineage holdout metadata fell to 0.297 average precision against
0.645 for the full stack. The random split measures a mixture, and a large part of what it measures
is available without looking at a gene.

\begin{figure}[t]\centering
\figbox{Figure 1. Study design}{6.5cm}{Left: the training collection (6,962 \Ecoli{}) and its
splits (random, phylogroup-stratified, ST holdout, ST-grouped 5-fold, country, time). Middle:
feature blocks (SNP blocks, presence/absence, per-gene statistics, determinant calls, metadata,
assembly statistics) and models (TriLens stack, its five classical learners, XGBoost and L1 on
determinant calls, determinant rule, assembly arm). Right: the external cohort (10,232 NCBI
isolates with MIC labels), the Giessen collection and the second species
(\textit{Campylobacter}, 4,426). Arrows mark which model is trained where and which is applied
frozen.}{schematic; no data file}
\caption{\textbf{Study design.} \todo{caption}}\label{fig:design}
\end{figure}

\begin{figure}[t]\centering
\figbox{Figure 2. The non-genomic floor}{6cm}{Grouped bars, mean over seven drugs, one group per
split (random, ST holdout, UK$\rightarrow$Norway, Norway$\rightarrow$UK, Norway time): submitting
study only, epidemiological fields only, assembly statistics only, all metadata, full genomic
stack. Panel A in average precision, panel B in \nap{}.}{results/diagnostics/confounder\_floor.csv}
\caption{\textbf{How much is predictable without reading a gene.} \todo{caption}}\label{fig:floor}
\end{figure}

\subsection*{Where the score comes from: components and baselines}

Removing one part of the stack at a time (Table~\ref{tab:ablation}) separates two claims that are
easy to conflate. The three branches averaged with equal weight reach 0.780; adding the gate brings
0.796; adding the five classical learners through the meta-learner brings 0.850. No part is dead
weight, and every drug improves at both steps, but the parts do not contribute equally: classical
stacking adds $+0.055$, more than three times the gate's $+0.016$. The model is a hybrid in which
the classical arm carries most of the load, and the classical arm on its own, stacked without the
neural branch, already reaches 0.846.

\begin{table}[t]\centering\small
\caption{\textbf{Component ablation on the random split}, mean average precision over the seven
drugs, with the gain each added part contributes over the stage before it.}
\label{tab:ablation}
\begin{tabular}{lrr}
\toprule
Stage & Mean AP & Gain \\
\midrule
Three branches, equally averaged & 0.780 & --- \\
$+$ mixture-of-experts gate (the neural arm alone) & 0.796 & $+0.016$ \\
Five classical learners, stacked, without the neural arm & 0.846 & --- \\
$+$ stacking the neural arm with them (full model) & 0.850 & $+0.055$ \\
\bottomrule
\end{tabular}
\end{table}

Against models from outside this work the lead is clear, and against its own strongest component it
is not (Table~\ref{tab:baselines}). The phylogroup label alone scores 0.291, against a mean
resistance frequency of 0.250 across the seven drugs: coarse lineage membership buys about 0.04
over chance and no more. TabNet \cite{arik2021tabnet}, built for tabular inputs, reaches 0.667 on
the random split, behind every one of the five classical learners, and falls to 0.341 on the
lineage holdout, where the weakest classical learner still reaches 0.440. This is the pattern that
tabular benchmarks outside this field report \cite{grinsztajn2022trees,shwartzziv2022tabular}. The
PCA model this project began with reaches 0.803. The full stack leads the best of these by 0.047,
TabNet by 0.183 and the null model by 0.559. Its margin over its own strongest component, XGBoost
on $\chi^2$-filtered SNPs, is $+0.008$ on the random split and $+0.023$ on the holdout --- inside
the interval widths above, and of a size a change of seed could produce. The defensible claim is
the weaker one: nothing simpler beat it, not that it beat everything simpler.

\begin{table}[t]\centering\small
\caption{\textbf{Baselines and components}, mean average precision over the seven drugs. The first
block is methods from outside this work; the second is the stack's own base learners, so those rows
are an ablation and not an external comparison. The phylogroup-only and PCA baselines were not
run on the lineage holdout.}
\label{tab:baselines}
\begin{tabular}{lrr}
\toprule
Model & Random split & Lineage holdout \\
\midrule
\multicolumn{3}{l}{\textit{External baselines}} \\
Phylogroup label only & 0.291 & not run \\
TabNet & 0.667 & 0.341 \\
XGBoost on 1,024 PCA components & 0.803 & not run \\
\midrule
\multicolumn{3}{l}{\textit{Base learners of the stack}} \\
Random forest, gene features & 0.781 & 0.493 \\
Logistic regression, gene features & 0.792 & 0.440 \\
LightGBM, gene features & 0.833 & 0.552 \\
XGBoost, gene features & 0.838 & 0.563 \\
XGBoost, $\chi^2$-filtered SNPs & 0.842 & 0.622 \\
Five classical learners, stacked & 0.846 & 0.622 \\
\midrule
Full stack & 0.850 & 0.645 \\
\bottomrule
\end{tabular}
\end{table}

\subsection*{Holding out whole lineages, with the prevalence effect removed}

The phylogroup-stratified split does not test generalisation, and it is worth seeing what it
reports: 0.856 against the random split's 0.850, four drugs above the random split and none more
than 0.022 below. Reshuffling strains inside eight large clusters leaves close relatives on both
sides, so no drop is possible. Under the sequence-type holdout, where no test strain shares a type
with any training strain and no test strain has a near-duplicate in training, the mean falls to
0.645.

Part of that fall is arithmetic. Forcing the four largest clones into training lowers the test
prevalence of every drug, by 2.5 to 12.1 percentage points; gentamicin's holdout test partition
contains 38 resistant strains, a prevalence of 4.2\% against 14.6\% in training. By the floor
described in Methods, a test partition with fewer resistant strains therefore scores lower even
when the model ranks strains identically. We measured the size of that effect
directly: for each drug we repeatedly subsampled susceptible strains in the holdout test partition
to the training prevalence and recomputed average precision on the same saved predictions, over
2,000 rounds, discarding no resistant strain and retraining nothing (Table~\ref{tab:previnv}).

\begin{table*}[t]\centering\footnotesize
\setlength{\tabcolsep}{4pt}
\caption{\textbf{The lineage holdout with the prevalence effect removed.} The prevalence-matched
column resamples susceptible strains in the holdout test partition to the training prevalence, over
2,000 rounds, on the same saved predictions. AUROC, whose no-skill value is 0.5 whatever the
prevalence, is given for comparison. The last column is the gap against the random split once
prevalence is matched.}
\label{tab:previnv}
\begin{tabular}{lrrlrr}
\toprule
Drug & Random split & Holdout, as measured & Holdout, prevalence-matched & AUROC & Adjusted gap \\
\midrule
AMP & 0.943 & 0.907 & 0.930 [0.920, 0.942] & 0.932 & $-0.014$ \\
CIP & 0.975 & 0.918 & 0.949 [0.932, 0.965] & 0.972 & $-0.026$ \\
SXT & 0.842 & 0.801 & 0.818 [0.806, 0.837] & 0.907 & $-0.025$ \\
\midrule
CTX & 0.904 & 0.647 & 0.752 [0.717, 0.788] & 0.885 & $-0.152$ \\
CAZ & 0.823 & 0.486 & 0.642 [0.597, 0.681] & 0.780 & $-0.181$ \\
AMC & 0.754 & 0.503 & 0.565 [0.545, 0.591] & 0.691 & $-0.189$ \\
GEN & 0.709 & 0.256 & 0.470 [0.407, 0.524] & 0.726 & $-0.239$ \\
\midrule
Mean & 0.850 & 0.645 & 0.732 & 0.842 & $-0.118$ \\
\bottomrule
\end{tabular}
\end{table*}

About two-fifths of the headline drop is a prevalence effect: the mean falls from 0.850 to 0.732
rather than to 0.645, a true gap of 0.118 instead of 0.205. Gentamicin recovers most, from 0.256 to
0.470, which its AUROC of 0.726 already implied: the model had not stopped ranking resistant strains
above susceptible ones, it was ranking them in a partition where one strain in twenty-four is
resistant.

Correcting for prevalence does not remove the finding; it sharpens it. The seven drugs fall into
two groups with nothing between them. Ampicillin, ciprofloxacin and
trimethoprim--sulfamethoxazole lose 0.014 to 0.026, which is nothing; cefotaxime, ceftazidime,
amoxicillin--clavulanic acid and gentamicin lose 0.152 to 0.239 (Figure~\ref{fig:twogroups}). A
model that had merely met a harder test would lose roughly evenly across the panel.
Amoxicillin--clavulanic acid is the clearest case: its prevalence shifts by only 5.9 points, so it
recovers least of the four hard drugs ($+0.062$ against gentamicin's $+0.214$) and has the lowest
AUROC of all seven at 0.691, leaving the least of its loss attributable to prevalence at all.

Two caveats attach to the table. Gentamicin's matched interval is wide, [0.407, 0.524], resting on
38 resistant strains, so its direction is clear and its value is not. And the matching target is the
training prevalence, which is slightly above the random-split test prevalence, which makes the
comparison a little generous to the holdout; the true gap may be marginally larger than 0.118.

\begin{figure}[t]\centering
\figbox{Figure 3. Two groups of drugs}{6.5cm}{Per drug, three bars: random split, lineage holdout
as measured, and holdout with prevalence matched to training. Drugs ordered by the
prevalence-corrected gap, printed above each group. Dashed line separates the three drugs losing
0.014--0.026 from the four losing 0.152--0.239.}{results/diagnostics/matched\_prevalence.csv;
results/pillar\_mlst\_final\_comparison/three\_way\_comparison\_FINAL.csv}
\caption{\textbf{The generalisation gap splits the panel in two.} \todo{caption}}\label{fig:twogroups}
\end{figure}

\subsection*{What the holdout costs clinically}

Within-class error rates are unaffected by the prevalence shift, because they are ratios inside one
class, so they need no correction before they can be read. At the fixed 0.5 threshold the mean very
major error rate rises from 0.241 on the random split to 0.559 on the holdout
(Table~\ref{tab:holdoutclin}). For three drugs the model has effectively stopped predicting
resistance at all: it identifies 6 of 71 cefotaxime-resistant strains, 7 of 43
ceftazidime-resistant strains and 5 of 38 gentamicin-resistant strains, with a major error rate of
0.001 or 0.002 in each case. A classifier that almost never calls resistance almost never calls it
wrongly; that pairing is the signature of collapse onto the majority class.

\begin{table*}[t]\centering\footnotesize
\setlength{\tabcolsep}{4pt}
\caption{\textbf{Clinical error rates under the lineage holdout} against the random split, at a
fixed threshold of 0.5. VME, very major error; ME, major error; CA, categorical agreement.}
\label{tab:holdoutclin}
\begin{tabular}{lrrrrrrr}
\toprule
Drug & Resistant in test & VME random & VME holdout & Increase & ME holdout & CA random & CA holdout \\
\midrule
CIP & 132 & 0.071 & 0.114 & $+0.043$ & 0.008 & 0.973 & 0.977 \\
AMP & 300 & 0.107 & 0.177 & $+0.070$ & 0.097 & 0.897 & 0.869 \\
SXT & 142 & 0.218 & 0.317 & $+0.099$ & 0.069 & 0.873 & 0.863 \\
\midrule
AMC & 145 & 0.394 & 0.683 & $+0.289$ & 0.073 & 0.805 & 0.763 \\
CAZ & 43 & 0.291 & 0.837 & $+0.546$ & 0.001 & 0.934 & 0.959 \\
GEN & 38 & 0.415 & 0.868 & $+0.453$ & 0.001 & 0.936 & 0.963 \\
CTX & 71 & 0.190 & 0.915 & $+0.725$ & 0.001 & 0.953 & 0.918 \\
\midrule
Mean & --- & 0.241 & 0.559 & $+0.318$ & --- & --- & --- \\
\bottomrule
\end{tabular}
\end{table*}

The categorical agreement columns are why this paper does not report accuracy. Three drugs agree
with the laboratory as well or better on the holdout than on the random split while their very major
error rates double or quadruple: gentamicin rises from 0.936 to 0.963 agreement while missing 33 of
38 resistant strains, ceftazidime from 0.934 to 0.959 while missing 36 of 43, and cefotaxime loses
0.035 of agreement while its very major error rate goes from 0.190 to 0.915. By the measure a
machine-learning paper usually reports, the model looks as good on unseen lineages; by the measure a
laboratory would use, it has stopped working.

This is the third analysis to produce the same two groups, on the same predictions, and the
grouping is stable even though the order within the second group is not: the three easy drugs gain
0.043--0.099 in very major error rate and the four hard ones 0.289--0.725.

The shortfall is not a matter of where the cut sits. Tuning a per-drug threshold on the validation
partition to maximise $F_2$ lowered the mean very major error rate on the random split from 0.241 to
0.133 while raising the mean major error rate from 0.055 to 0.141; constraining the major error rate
to the 3\% benchmark instead pushed the very major error rate to 0.308. No drug met both limits at
any threshold tested, because moving the cut trades one error against the other. What limits the
model is class separation, not calibration.

\subsection*{Transfer of the recipe to the Giessen collection}

The two collections share no feature vector, so this is a comparison of the method and not of a
model: the classical stacking half of the pipeline was retrained on Giessen with the same protocol
and hyperparameters, and the neural branches, which need the pan-genome feature shapes, are absent
from both sides of the comparison. On the four shared drugs it reaches a mean of 0.769 against
0.849 for the same half on the training collection, 89.6\% of the per-drug ratio
(Table~\ref{tab:giessen}).

\begin{table}[t]\centering\small
\caption{\textbf{The classical stacking half of the pipeline, retrained on an independent
collection.} Neither column is the full stack. Retention is inflated by Giessen's higher
prevalence, as the text explains.}
\label{tab:giessen}
\begin{tabular}{lrrrr}
\toprule
Drug & Training collection & Giessen & Prevalence (train / Giessen) & Retention \\
\midrule
CIP & 0.976 & 0.988 & 0.228 / 0.452 & 101.2\% \\
CAZ & 0.820 & 0.763 & 0.131 / 0.341 & 93.1\% \\
CTX & 0.895 & 0.786 & 0.152 / 0.443 & 87.8\% \\
GEN & 0.707 & 0.539 & 0.118 / 0.232 & 76.3\% \\
\midrule
Mean & 0.849 & 0.769 & --- & 89.6\% \\
\bottomrule
\end{tabular}
\end{table}

That retention figure cannot be read as preserved accuracy, for the reason the prevalence correction
above established. Every Giessen drug has two to three times the resistance prevalence of its
training counterpart, so every no-skill baseline is two to three times higher before the model is
credited with anything. Ciprofloxacin's apparent 101\% is the clearest case: the task got slightly
easier, not the model slightly better. Prevalence-matching cannot fix it here, because matching
would require resampling both collections to a common reference and the 81-strain Giessen test
partition will not survive that. The claim this table supports is the weaker one, that the recipe
reconstructs on a genomically different collection rather than falling apart, and that its design
choices are not fitted to one dataset. The interesting part is the ordering: ciprofloxacin retains
most and gentamicin least, the same ordering as under the lineage holdout, reached by a different
route.
